% !TeX root = ../../main.tex
\subsection{Proof of Theorem 1.1}
\label{subsec:s04:04}

It remains to combine the finite coefficient calculations with the
uniform estimates for the whole graph sum.

\begin{prfn}{Theorem~\ref{thm:lambda-one}}
Apply Proposition~\ref{prp:s03:graph-tail} with $M=3$ at $u=1$.
The absolute sum of graphs without a large vertex or with $q\geq3$
is $O(F_g g^{-3})$. There are only finitely many remaining core
types. First discard their legal loop counts $p>\log g$ by
Lemma~\ref{lmm:s03:loop-tail}. This tail is smaller than every fixed
inverse power after normalization, including when the large vertex's
genus is small.

For a remaining core of complexity $q\leq2$, apply
Proposition~\ref{prp:s03:core-expansion} with $K=2-q$.
All its intermediate volume types meet the uniform input thresholds
for $p\leq\log g$. There are constants $C>0$ and a fixed integer
$s\geq0$, depending on the core and expansion order, for which the
remainder after division by $V_{g,0}$ is bounded by
\[
 g^{-3}\frac{C^{p+1}(1+p)^s}{p!}.
\]
This bound sums to $O(g^{-3})$ over $p$.
The same factorial bound on the coefficient functions permits
extending each retained coefficient sum to all $p\geq0$, with a
further superpolynomially small tail. This is the extension used in
the signed moments above; no large-genus expansion has been applied
to an illegal or small-genus vertex.

Multiplication by~\eqref{eq:s03:fixed-n-volume} at $n=0$ through
order two converts the relative core expansions to the absolute
groups~\eqref{eq:s04:absolute-groups}. It preserves their summable
remainder bound and supplies the leading WP factor $1/\sqrt\pi$.
Consequently
\begin{align*}
 \frac{Z_g(1)}{F_g}=\frac{e^{S(1)}}{\sqrt\pi}
 \left[Q_{0,0}+\frac{Q_{0,1}+Q_{1,0}}g\right.\qquad\qquad\\
 \left.{}+\frac{Q_{0,2}+Q_{1,1}+Q_{2,0}}{g^2}
                  +O(g^{-3})\right].
\end{align*}
Propositions~\ref{prp:s04:first-correction}
and~\ref{prp:s04:second-correction} give the two coefficient sums.
Finally~\eqref{eq:s03:normalization-links} restores
$a_g/N_g=Z_g(1)/F_g$, and~\eqref{eq:s03:cin-identity} gives
$S(1)=\Cin(2\pi)$. Thus
\[
 a_g=\frac{e^{\Cin(2\pi)}}{\sqrt\pi}N_g
 \left[1+\frac9{8g}
     +\left(\frac{257}{128}+\frac1{8\pi^2}\right)g^{-2}
     +O(g^{-3})\right],
\]
which is~\eqref{eq:111}.
\end{prfn}
